Com valor inicial $P_0=3$ bilhoes em 1960,
\[
P(t)=3(1+r)^{t-1960},
\qquad r\in\{0.01,0.02,0.03\}.
\]

Em papel semi-logaritmico de base 10, tomamos logaritmo:
\[
\log_{10}P(t)=\log_{10}3+(t-1960)\log_{10}(1+r).
\]
Portanto, as tres curvas aparecem como retas.

As inclinacoes sao:
\[
r=1\%: \ \log_{10}(1.01)\approx 0.00432\ \text{por ano},
\]
\[
r=2\%: \ \log_{10}(1.02)\approx 0.00860\ \text{por ano},
\]
\[
r=3\%: \ \log_{10}(1.03)\approx 0.01284\ \text{por ano}.
\]

A interceptacao comum em $t=1960$ e
\[
\log_{10}3\approx 0.4771.
\]

Assim, em papel semi-log as tres curvas sao retas de mesma ordenada inicial e inclinacoes crescentes; a curva de $3\%$ e a mais inclinada, a de $1\%$ a menos inclinada, e todas representam crescimento exponencial puro.